% =====================================================================
%  CHAPTER 9: 儿童少年传染病（对应大纲第九章）
%  大纲要求：儿童少年传染病流行特征 [双横线/掌握]
%            儿童少年传染病预防控制 [双横线/掌握]
% =====================================================================
\chapter{儿童少年传染病}
\setcounter{chapter}{9}
\chapterintro{
掌握传染病的定义及常见类型；掌握儿童少年传染病的流行特征（疾病分布、人群分布、地区分布、时间分布）；掌握学校传染病流行过程特点（传染源集散地、传播特征与学生生活特点、儿童少年易感性，包括城市流动儿童和农村留守儿童）；掌握学校传染病预防控制措施（控制传染源、切断传播途径、保护易感人群、学校传染病管理制度）。
}

% ----- 10.1 流行特征 -----
\section{儿童少年传染病流行特征}

\subsection{传染病的定义与分类}

\begin{keybox}
\textbf{一、传染病（infectious disease）：}由病原体（细菌、病毒、立克次体、螺旋体等）引起的，能在人与人、人与动物或动物与动物间相互传染的疾病，是许多疾病的总称。
\begin{enumerate}[leftmargin=*]
    \item \imp{病毒性传染病}：流感（季节性流感）、腮腺炎、手足口病、水痘、麻疹、轮状病毒肠炎（秋季腹泻）。
    \item \imp{细菌性传染病}：猩红热、百日咳、细菌性痢疾。
\end{enumerate}
\end{keybox}

\subsection{儿童少年传染病流行状况}

\begin{keybox}
\textbf{一、疾病分布}
\begin{enumerate}[leftmargin=*]
    \item \imp{疾病种类分布广泛}：儿童少年中报告的传染病种类占我国法定传染病种类的80\%，其中以呼吸道疾病最多。
    \item \imp{呼吸道传染病仍是学校传染病预防的重点}：呼吸道传染病以小学生最多，高中生和大学生次之，初中生最少；肠道传染病也是小学生最多，其余学段基本持平。
    \item \imp{高发疾病相对集中}：季节性流感、流行性腮腺炎、手足口病、风疹等。
\end{enumerate}

\textbf{二、人群分布}
\begin{enumerate}[leftmargin=*]
    \item \imp{学段}：不同传播途径和高发传染病在各学段中分布不同。
    \item \imp{性别}：学生中甲乙丙类传染病报告发病数各年龄组男生均高于女生。
    \item \imp{年龄}：学生甲乙类传染病发病数的年龄分布呈"马鞍形"。
\end{enumerate}

\textbf{三、地区分布}：不同省份间传染病报告发病数波动范围大。原因：
\begin{enumerate}[leftmargin=*]
    \item 地区经济发展水平差异；
    \item 人口密度与流动情况；
    \item 地理气候条件（南方省份气候温暖湿润，适合蚊虫等病媒生物滋生繁殖）。
\end{enumerate}

\textbf{四、时间分布}：总体呈双峰分布，全年呈现2个发病高峰，主要集中在3-6月和10-12月。
\end{keybox}

\subsection{学校传染病流行过程特点}

\begin{keybox}
\textbf{一、传染源的集散地}
\begin{enumerate}[leftmargin=*]
    \item \imp{人群庞杂}：学校群体年龄结构包括儿童、青少年和成年人。
    \item \imp{流动性强}：日常教学与社交活动导致大量跨区域、跨群体人员接触，为病原体传播创造接触网络。
    \item \imp{健康携带者众多}：许多传染病除患者外，还有数倍甚至数十倍的健康携带者，使传染源数量、强度显著增加。
\end{enumerate}

\textbf{二、传播特征与学生生活特点有关}
\begin{enumerate}[leftmargin=*]
    \item \imp{群体性生活方式}：集中住宿、共用餐区和统一作息显著增加接触密度，呼吸道飞沫/接触传播风险增加。
    \item \imp{室内微小环境中细菌聚集}：空气流通受限的室内环境中，病原体气溶胶浓度比开放空间高很多。
    \item \imp{公共设施使用频率高}：门把手、座椅、图书馆等。
    \item \imp{学习制度}：固定课表、集中授课等教学安排强制维持群体接触状态，日均暴露时长>6h，远超社区接触强度。
\end{enumerate}

\textbf{三、儿童少年易感性}
\begin{enumerate}[leftmargin=*]
    \item \imp{年龄}：年龄越小机体抵抗力越低，免疫系统作为身体抵御病原体的重要防线，婴幼儿胸腺、脾脏等免疫器官发育不完全，免疫细胞数量和活性也相对较低，使他们更易受各种传染病侵袭，如呼吸道合胞病毒感染。
    \item \imp{特异性免疫能力}：①机体特异性免疫功能尚不完善；②群体内基础免疫水平不一。
    \item \imp{卫生习惯不良}。
    \item \imp{传染病相关知识掌握不足}：不了解传染病传播途径和预防。
    \item \imp{特殊儿童群体免疫接种不全}：
    \begin{itemize}
        \item \imp{城市流动儿童}：由于所在地经常变动，他们的预防接种管理往往不规范或有所遗漏，经常发生迟种、不种、接种间隔时间不规范的情况，使流动儿童发生传染病的几率比普通儿童高很多。影响流动儿童疫苗接种情况的主要因素：\\
        A. 人口流动性：在流入地居住时间，异地来往频繁程度；\\
        B. 年龄：年龄越小的儿童，能完成规划免疫疫苗人数越多；\\
        C. 出生地和接生场所；\\
        D. 监护人个人素质（尤其是母亲文化程度）：文化程度越高、有职业，免疫接种知识掌握程度越高，儿童越能按时完成疫苗接种；\\
        E. 预防接种服务供给利用情况：不同城市预防接种服务资源分布不均衡。
        \item \imp{农村留守儿童}：\\
        A. 传染病罹患率高：健康知识缺乏、卫生习惯较差、营养摄入不足、医疗资源匮乏；\\
        B. 疫苗接种率低：按期接种、全程接种坚持性差，超期接种、免疫空白等现象多发；\\
        C. 监护人意识淡薄，交通不便，信息沟通不畅。
    \end{itemize}
\end{enumerate}

\textbf{四、传染病对儿童少年健康的影响}
\begin{enumerate}[leftmargin=*]
    \item \imp{生理健康}：(1)儿童少年罹患传染病后通常会出现发热、头痛、出疹、食欲降低等；(2)未及时治疗会引发严重并发症。
    \item \imp{成年期生活质量}。
    \item \imp{疾病负担}。
\end{enumerate}
\end{keybox}

% ----- 10.2 传染病预防控制 -----
\section{儿童少年传染病预防控制}

\subsection{控制传染源}

\begin{keybox}
\textbf{一、管理传染源}
\begin{enumerate}[leftmargin=*]
    \item \imp{早发现}：坚持晨检制度并保持经常化，是早期发现传染病病例的重要保障。通过每日入校前健康检查，及时发现发热、皮疹、腹泻等症状的学生。
    \item \imp{早报告}：发现传染病病例或疑似病例后，应在第一时间内向疾病预防控制机构和教育行政部门报告，以便采取正确办法控制疫情。
    \item \imp{早隔离}：对确诊病例、疑似病例和早期症状者，应及时采取隔离措施。隔离是保护易感人群的必要措施，防止病原体在校园内进一步传播。
    \item \imp{早治疗}：对确诊病例、疑似病例和可疑病例的早期症状者，应根据疾病特点，及时将患者送定点医院隔离治疗或在家隔离治疗，防止病情加重和传播扩散。
\end{enumerate}

\textbf{二、"四早"措施}：传染病预防控制的"四早"措施即\imp{早发现、早隔离、早报告和早治疗}。

\textbf{三、疫情报告制度}：当出现以下情况之一时，应在24小时内报出相关信息：
\begin{enumerate}[leftmargin=*]
    \item 在同一宿舍或者同一班级，1天内有3例或者连续3天内有多名学生（5例以上）患病，并有相似症状（如发热、皮疹、腹泻、呕吐、黄疸等）或者有共同用餐、饮水史；
    \item 个别学生出现不明原因的高热、呼吸急促或剧烈呕吐、腹泻等症状；
    \item 学校发生群体性不明原因疾病或者其他突发公共卫生事件。
\end{enumerate}
当发生大面积疫情时必须实行"零报告"和日报告制度。
\end{keybox}

\subsection{切断传播途径}

\begin{keybox}
\textbf{一、针对不同传播途径的切断措施}
\begin{enumerate}[leftmargin=*]
    \item \imp{呼吸道传染病}：加强室内通风换气，保持空气流通；定期进行空气消毒；流行季节减少大型集会和集体活动；提倡佩戴口罩。
    \item \imp{肠道传染病}：加强饮水和食品安全管理；对粪便进行无害化处理；消毒被污染物品和周围环境；做好食堂、厕所等重点场所的卫生管理。
    \item \imp{接触传播}：对门把手、桌椅、楼梯扶手、图书馆书籍等公共设施进行定期消毒；培养学生勤洗手的卫生习惯。
    \item \imp{虫媒传染病}：消灭虫媒传播的中间宿主，如灭蚊、灭蝇、灭鼠等；改善校园环境卫生，清除蚊虫滋生地。
    \item \imp{血液/体液传播}：加强健康教育，杜绝吸毒和共用注射器；做好医疗器械的消毒管理工作。
\end{enumerate}

\textbf{二、学校环境卫生管理}
\begin{enumerate}[leftmargin=*]
    \item 保持教室、宿舍、食堂等场所的清洁卫生和通风良好。
    \item 落实公共区域（门把手、楼梯扶手、卫生间、水龙头等）的日常消毒制度。
    \item 加强学校饮用水和食品安全管理，防止水源性和食源性传染病暴发。
    \item 做好垃圾和污物的无害化处理，消除病媒生物滋生环境。
\end{enumerate}
\end{keybox}

\subsection{保护易感人群}

\begin{keybox}
\textbf{一、计划免疫接种}
\begin{enumerate}[leftmargin=*]
    \item \imp{常规免疫接种}：按照国家免疫规划程序，组织和督促学生按时完成各类疫苗的接种，是预防儿童少年传染病发生的最有效方法。
    \item \imp{应急接种}：在传染病暴发流行时，根据疾病预防控制机构的指导，对易感人群开展应急接种。
    \item \imp{查验预防接种证}：入学入托时查验预防接种证，对未完成接种的学生督促补种。
\end{enumerate}

\textbf{二、健康教育}
\begin{enumerate}[leftmargin=*]
    \item 开展形式多样的传染病预防性健康教育，普及传染病传播途径和预防知识。
    \item 根据传染病流行特点（季节性、地区性）开展针对性健康教育。
    \item 培养学生良好的个人卫生习惯：勤洗手、不随地吐痰、不共用个人物品等。
    \item 健康教育对象不仅包括学生，还应覆盖教师、家长和后勤人员。
\end{enumerate}

\textbf{三、增强体质}
\begin{enumerate}[leftmargin=*]
    \item 保证学生充足的睡眠和合理的营养，增强机体非特异性免疫力。
    \item 鼓励学生积极参加体育锻炼，提高身体素质。
    \item 合理安排学习与休息时间，避免过度疲劳导致抵抗力下降。
    \item 流行季节注意防寒保暖，减少感染机会。
\end{enumerate}

\textbf{四、特殊群体保护}
\begin{enumerate}[leftmargin=*]
    \item \imp{城市流动儿童}：健全组织管理网络，提高疫苗接种便捷性；加强监护人健康教育，提高其免疫接种知识和意识；确保流动儿童在流入地享有同等的预防接种服务。
    \item \imp{农村留守儿童}：加强健康教育力度，普及传染病预防知识；改善卫生习惯和营养状况；提高疫苗接种的可及性和按时接种率，消除免疫空白；建立健全监护人联系机制，确保信息沟通畅通。
\end{enumerate}
\end{keybox}

\subsection{学校传染病管理制度}

\begin{keybox}
\textbf{一、晨检制度}
每天早晨由班主任或校医对学生进行健康观察，重点检查有无发热、皮疹、咽痛、腹泻等症状。发现异常情况及时报告校医，进行初步排查和处置。晨检制度是"早发现"的重要保障，必须坚持经常化、制度化。

\textbf{二、因病缺勤追踪与登记制度}
对因病缺勤的学生，班主任应及时了解其患病情况和可能的病因，做好登记记录。发现传染病或疑似传染病病例时，应立即向校医和学校领导报告。校医负责追踪学生的诊断和治疗情况，判断是否构成疫情。

\textbf{三、传染病疫情报告制度}
学校应指定专人（校医或保健教师）作为学校传染病疫情报告人，负责传染病的监测、报告和登记工作。疫情报告人应接受专业培训，掌握传染病诊断标准和报告流程。发现法定传染病或聚集性病例时，应在规定时限（24小时内）向属地疾病预防控制机构和教育行政部门报告。

\textbf{四、复课证明查验制度}
患传染病的学生在隔离期满、症状消失后，须凭医疗机构出具的康复证明方可返校复课。校医负责查验复课证明，确认学生已无传染性后方可解除隔离。

\textbf{五、新生入学预防接种证查验制度}
在新生入学入托时，学校应查验其预防接种证。对未按规定完成免疫接种的学生，应督促其家长带领学生到预防接种单位进行补种，并将补种情况记录归档。

\textbf{六、学校卫生消毒制度}
建立教室、宿舍、食堂、图书馆等重点场所的定期消毒制度，明确消毒频次、消毒方法和责任人。传染病流行期间应增加消毒频次，并做好消毒记录。

\textbf{七、校医职责}
校医是学校卫生专业队伍中的重要组成部分，是国家学校卫生工作策略的贯彻者和执行者，是学生和教师健康的守护者。校医在学校传染病监测、预防控制和应急处置中发挥着重要作用，具体职责包括：(1)负责晨检、因病缺勤追踪和疫情报告工作；(2)组织开展传染病预防健康教育；(3)落实学校卫生消毒和传染病防控措施；(4)参与学校突发公共卫生事件的应急处置。

\textbf{八、建立学校公共卫生预警系统}
建立健全学校传染病监测预警机制，对晨检、因病缺勤、疫情报告等信息进行分析研判，及时发现传染病早期聚集性风险，实现传染病疫情的早发现、早预警、早处置。
\end{keybox}

% ===== 儿童少年常见传染病各论 =====
\section{儿童少年常见传染病各论}

\subsection{流行性感冒}

\begin{keybox}
\textbf{流行性感冒（influenza）：}由流感病毒引起的急性呼吸道传染病，以急起高热、全身酸痛、乏力，或伴轻度呼吸道症状为主要临床特点。潜伏期短，传染性强，传播迅速。

\textbf{流行特征：}流感病毒分甲、乙、丙三型，甲型流感威胁最大。流行特征表现为突然发生，迅速蔓延，发病率高和流行过程短。流行无明显季节性，以冬春季节为多。

\textbf{传播途径：}主要传染源是患者和隐性感染者。患者自潜伏期末到发病后5日内均可排出病毒，传染期约1周，以病初2～3日传染性最强。以飞沫传播为主。

\textbf{预防措施：}
\begin{enumerate}[leftmargin=*]
    \item 早期发现，及时报告、隔离和治疗患者。
    \item 接种疫苗：预防该疾病或疾病严重后果的最有效方法。
    \item 流行季节避免大型集会，注意休息，防止过劳。
    \item 注意体育锻炼和营养，加强非特异性免疫。
\end{enumerate}
\end{keybox}

\subsection{结核病}

\begin{keybox}
\textbf{结核病（tuberculosis, TB）：}由结核分枝杆菌引起的一种可防可治的疾病。肺结核是最为常见的一种结核病，由结核分枝杆菌在肺部感染引起，是对健康危害较大的慢性传染病。

\textbf{流行特征：}结核病是全球重大公共卫生问题。中国是结核病高负担国家之一。儿童少年结核感染率随年龄增长而升高。学校作为人群密集场所，易发生结核病聚集性疫情。近年来学生肺结核报告发病率呈波动趋势，部分省份学生病例数占全人群病例数的5\%-8\%。高中和大学阶段是学生结核病的高发期。

\textbf{传播途径：}主要传播途径为经空气传播。患者咳嗽排出而悬浮于空气飞沫中的结核菌被健康人吸入后即可引起感染。与患者共餐、共用碗筷或喝了被结核菌污染的生牛奶等，均有被感染的可能。

\textbf{预防措施：}
\begin{enumerate}[leftmargin=*]
    \item 积极开展"结核病防治"健康教育，普及卫生知识。
    \item 定期筛查。
    \item 积极治疗现患者，切断传播途径。
    \item 通过药物预防易感染者发病。
    \item 预防接种：通过接种卡介苗以获得对结核菌的特异性免疫力。
\end{enumerate}
\end{keybox}

\subsection{性传播疾病}

\begin{defbox}
\textbf{性传播疾病（sexually transmitted disease, STD）：}由性接触、类似性行为及间接接触所感染的一组传染性疾病。性病可由30多种性传播的病毒、细菌和寄生虫引起。
\end{defbox}

\begin{keybox}
\textbf{流行特征：}WHO报道每年全世界估计有3.4亿名15-24岁青少年感染可治愈的性传播疾病。青少年STD感染率在全球范围内呈上升趋势。我国青少年性传播疾病以梅毒、淋病和艾滋病为主要病种。HIV感染与结核病构成青春期青少年传染病疾病负担的重要原因。青少年由于缺乏预防知识、不善于寻求卫生服务、存在高危行为等因素，成为STD的易感人群。

\textbf{性病的传播方式（6种）：}
\begin{enumerate}[leftmargin=*]
    \item \imp{性接触传播}：性交行为是主要传播途径。
    \item \imp{非性行为的直接接触传播}：通过直接接触患者的病变部位或其分泌物所致。
    \item \imp{血源感染}：经静脉输入感染的血液或血液制品，及静脉注射毒品等。
    \item \imp{母婴传播}：母亲携带的病原体可经胎盘、产道等途径传染给胎儿或新生儿。
    \item \imp{医源性传播}：医务人员操作中因防护不当或使用未严格消毒的器械导致感染。
    \item \imp{日常生活接触传播}。
\end{enumerate}

\textbf{青少年主要性传播疾病：}WHO报道每年全世界估计有3.4亿名15～24岁青少年感染可治愈的性传播疾病。HIV感染与结核病构成青春期青少年传染病疾病负担的重要原因。

\textbf{梅毒（syphilis）：}由梅毒螺旋体引起的慢性、系统性性传播疾病。一期梅毒：入侵部位出现硬下疳，多数发生于外生殖器部位；二期梅毒：感染后7～10周，腹股沟淋巴结肿大；三期梅毒（晚期）：病期超过2年，典型表现为内脏梅毒，累及心血管、肝脏、神经等。

\textbf{淋病（gonorrhea）：}由淋病奈瑟菌引起的泌尿和生殖系统化脓性感染。潜伏期平均3～5天。男性：尿道口红肿、尿急、尿频、尿痛、尿道口流脓；女性：外阴瘙痒、白带增多、宫颈充血、触痛、脓性分泌物。

\textbf{艾滋病（AIDS）：}由人类免疫缺陷病毒（HIV）引起，HIV攻击和破坏人体免疫系统，导致继发性感染和癌症，最终导致死亡。

\textbf{性传播疾病的预防原则：}
\begin{enumerate}[leftmargin=*]
    \item 做好性病的监测工作
    \item 普及预防性病知识
    \item 重视学校健康教育
    \item 不歧视HIV感染者和AIDS患者
    \item 积极正规治疗
    \item 关注重点人群（流动青少年等）
\end{enumerate}
\end{keybox}

\subsection{流行性腮腺炎}

\begin{keybox}
\textbf{流行性腮腺炎（mumps）：}由腮腺炎病毒引起的急性传染病，通常引发腮腺疼痛红肿、发烧、头痛和肌肉酸痛。

\textbf{流行特征：}人群普遍易感，5～15岁儿童青少年为高危人群。发病季节呈双峰分布。

\textbf{传播途径：}早期患儿和隐性感染者是主要传染源。通过空气中飞沫经呼吸道传播。唾液污染的食物、餐具和其他用品也可传播。

\textbf{预防措施：}
\begin{enumerate}[leftmargin=*]
    \item 主动免疫：疫苗接种。
    \item 被动免疫：儿童在罹患该病后可获得较为持久的免疫力。
    \item 隔离患儿，早发现早隔离：腮腺肿胀消退1周后方可解除隔离，密切接触者应检疫21天。
    \item 切断传播途径。
\end{enumerate}
\end{keybox}

\subsection{手足口病}

\begin{keybox}
\textbf{手足口病（hand-foot-mouth disease, HFMD）：}由多种肠道病毒引起，以发热和手部、足部、口腔等部位出现皮疹或疱疹为特征的急性传染病。以柯萨奇病毒A16型和肠道病毒71型（EV71）最为常见。

\textbf{流行特征：}手足口病是我国丙类传染病中报告发病率最高的病种之一，5岁以下儿童占报告病例的90\%以上，3岁以下儿童重症和死亡风险最高。全年均有发病，南方地区5-7月和9-11月两个流行高峰，北方地区以6-8月夏季高峰为主。托幼机构是手足口病暴发的主要场所。近年来EV71疫苗的推广使重症和死亡病例大幅减少。

\textbf{传播途径：}病毒可通过直接接触感染者鼻咽分泌物、唾液、疱疹液或粪便传播。被粪便、疱疹液、分泌物污染的物品也可传播。感染者发病后一周内传染性最强，传染期可持续数周。

\textbf{预防措施（WHO指南）：}
\begin{enumerate}[leftmargin=*]
    \item 建立和加强疾病监测体系。
    \item 开展良好的卫生习惯和基础环境卫生教育活动。
    \item 在疾病暴发时，对幼儿园、日托机构、学校提供援助。
    \item 加强对卫生保健机构和社区的感染控制措施。
    \item 改善临床患儿管理服务。
    \item 交换并传播与预警、应急响应和疾病管理有关的信息。
    \item 建立国家层面的行政管理框架。
    \item 对预防管理措施进行监测和评估。
\end{enumerate}
\end{keybox}

\subsection{诺如病毒肠炎}

\begin{keybox}
\textbf{诺如病毒肠炎（norovirus enteritis）：}由诺如病毒引起的，以恶心、呕吐、腹痛和腹泻等为主要症状的胃肠道疾病。成人腹泻较突出，儿童则较多有呕吐。

\textbf{流行特征：}全年均可发生，每年10月到次年3月是诺如病毒感染高发季节。

\textbf{传播途径：}摄入被粪便或呕吐物污染的食物或水；接触患者粪便或呕吐物；吸入呕吐时产生的气溶胶；间接接触被粪便或呕吐物污染的物品和环境。

\textbf{预防措施：}
\begin{enumerate}[leftmargin=*]
    \item 养成良好的卫生习惯。
    \item 注意饮食卫生。
    \item 切断传播途径。
\end{enumerate}
\end{keybox}

\newpage
